 \chapter{化学信使与信号转导}\label{ch:b2:cell-signalling}

一个肾上腺素分子到达肝细胞的表面，
从不进入细胞。它与外侧的一种蛋白质结合，使其
改变形状；不到一秒钟，细胞便已生成一万个
\omterm{prop:b2:cell-signalling:gpcr}{第二信使}分子，激活了十万个
酶分子，并开始释放一百万个
葡萄糖分子——一条长达百万倍的放大链，由
一个停留在门外的信号启动。本章讨论细胞之间发送的
信使，以及细胞如何将表面信息转化为
其化学过程、基因或
形状的改变：受体、\omterm{prop:b2:cell-signalling:gpcr}{第二信使}、激酶的
级联，以及两种设计——神经的快速有线信号，
与\omterm{def:b2:cell-signalling:kinds}{激素}的缓慢广播式信号传递——二者的对比
贯穿本书的其余部分。

\section{信使}

\begin{definition}[化学信号传递的类型]\label{def:b2:cell-signalling:kinds}
\index{激素}\index{内分泌信号传递}\index{旁分泌信号传递}\index{神经递质}
细胞借分泌的分子向另一细胞传递信号，该分子是
\emph{信使}或\emph{配体}，它与靶细胞上的\emph{受体}结合，受体位于
靶细胞表面或内部。按作用范围分：\emph{内分泌}信号传递，其中
\emph{激素}由腺体释放入\omterm{def:b2:blood-circulation:blood}{血液}，到达
全身每个细胞，仅作用于带有相应受体的细胞，
作用时间为数分钟至数小时（胰岛素、皮质醇、肾上腺素、甲状腺素）；
\emph{旁分泌}信号传递，信使扩散至
邻近细胞，扩散范围在一毫米之内，途中即被降解
（生长因子、组胺、前列腺素，以及
\cref{ch:b2:vertebrate-development}中的\omterm{thm:b2:vertebrate-development:gradient}{形态发生素}）；\emph{自分泌}，即细胞向自身发信号；
以及\emph{突触信号}，其中神经元
跨越\qty{20}{nm}的间隙释放\emph{神经递质}，
作用于单个靶细胞，仅需一毫秒（见\cref{ch:b2:neurons-synapses}）。
按化学性质分：\emph{水溶性}信使——氨基酸
衍生物（如肾上腺素）、肽及蛋白质（如
胰岛素）——不能穿过膜，只能作用于表面受体；
\emph{脂溶性}信使——类固醇（皮质醇、
雌二醇、睾酮）、甲状腺素、视黄酸、一氧化氮——
可穿过膜，作用于细胞内的受体。内分泌系统与神经系统
是由 $10^{13}$ 个细胞组成的机体协调自身的两种方式：
一种缓慢地向全身广播，另一种快速地
定向传给某一个细胞。
\end{definition}

\begin{omfigure}
\begin{tikzpicture}[every node/.style={font=\scriptsize}, >=stealth]
  % endocrine
  \begin{scope}
    \draw[thick, fill=omProp!25] (0,1.6) circle (0.35); \node[above] at (0,2.0) {腺体};
    \draw[very thick, omThm!60] (-0.5,0.6) -- (4.5,0.6); \draw[very thick, omThm!60] (-0.5,0.2) -- (4.5,0.2); \node[omThm!70!black] at (2,0.4) {血液};
    \foreach \x in {0.6,1.5,2.4,3.3} \fill[omProp] (\x,0.4) circle (0.05);
    \draw[->, omProp] (0,1.25) -- (0,0.65);
    \draw[thick, fill=omDef!20] (3.8,-0.8) circle (0.35); \node[below, align=center] at (3.8,-1.2) {靶细胞 \\有受体};
    \draw[->, omProp] (3.8,0.15) -- (3.8,-0.4);
    \draw[thick, fill=gray!20] (1.0,-0.8) circle (0.35); \node[below, align=center] at (1.0,-1.2) {无受体：\\不响应};
    \node[align=center] at (2.2,-2.4) {内分泌：激素入血，\\遍及全身，数分钟至数小时};
  \end{scope}
  % paracrine
  \begin{scope}[xshift=6.6cm]
    \draw[thick, fill=omProp!25] (0.8,0.6) circle (0.35);
    \foreach \a in {0,45,...,315} \fill[omProp] (0.8,0.6) ++(\a:0.6) circle (0.04);
    \draw[thick, fill=omDef!20] (1.9,0.9) circle (0.3); \draw[thick, fill=omDef!20] (1.7,-0.2) circle (0.3); \draw[thick, fill=omDef!20] (-0.2,0.0) circle (0.3);
    \node[align=center] at (0.8,-2.4) {旁分泌：扩散至 \\邻近细胞，\\一毫米内即降解};
  \end{scope}
  % synaptic
  \begin{scope}[xshift=11.6cm]
    \draw[very thick, omDef] (0,1.8) -- (0,0.5); \draw[thick, fill=omDef!25] (0,0.3) circle (0.25);
    \draw[thick, fill=omExo!30] (0,-0.5) circle (0.3);
    \foreach \x in {-0.08,0,0.08} \fill[omProp] (\x,-0.05) circle (0.03);
    \node[align=center] at (0.2,-2.4) {突触信号：单个靶细胞，\\\qty{20}{nm}，\\仅需一毫秒};
  \end{scope}
\end{tikzpicture}

{\small 化学信号的三种作用范围：\omterm{def:b2:cell-signalling:kinds}{激素}经
\omterm{def:b2:blood-circulation:blood}{血液}送达所有带相应受体的细胞；旁分泌因子
到达邻近细胞；\omterm{def:b2:cell-signalling:kinds}{神经递质}经
突触送达单个细胞。}
\end{omfigure}

\begin{theorem}[受体占有率]\label{thm:b2:cell-signalling:occupancy}
\index{解离常数}\index{受体占有率}
受体 $R$ 与其配体 $L$ 可逆结合，$R + L
\rightleftharpoons RL$，其\emph{\omterm{thm:b2:cell-signalling:occupancy}{解离常数}}为 $K_d =
[R][L]/[RL]$。平衡时，被占据的受体比例为
\[
\theta = \frac{[RL]}{[R]_{\text{总}}} = \frac{[L]}{K_d + [L]},
\]
与酶的饱和作用为同一双曲线。$K_d$ 是
使半数受体被占据的浓度，它度量
\emph{亲和力}：数值越低，结合越紧密。\omterm{def:b2:cell-signalling:kinds}{激素}受体的
$K_d$ 处于纳摩尔至皮摩尔量级，与\omterm{def:b2:cell-signalling:kinds}{激素}
循环时的浓度相匹配（$10^{-12}$ 至
$10^{-9}$ mol/L）：当 $[L] = K_d/10$ 时，十分之一被占据；当
$10K_d$ 时十分之九被占据。细胞无法响应
没有相应受体的\omterm{def:b2:cell-signalling:kinds}{激素}，它通过改变
受体的数目来调节敏感性，或者——如\cref{ch:b2:reproductive-hormones}中所述——
在\omterm{def:b2:cell-signalling:kinds}{激素}持续存在时使受体撤出。
\end{theorem}
\begin{proof}
由 $[R] = [R]_{\text{总}} - [RL]$，$K_d = ([R]_{\text{总}} -
[RL])[L]/[RL]$，故 $[RL](K_d + [L]) = [R]_{\text{总}}[L]$，即为
所求公式。它就是本科一年级分册中的米氏
表达式去掉催化步骤后的形式。
\end{proof}

\begin{omfigure}
\begin{tikzpicture}
\begin{axis}[
  width=0.68\linewidth, height=5.2cm,
  axis x line=bottom, axis y line=left,
  xmode=log, xmin=0.01, xmax=100, ymin=0, ymax=1.1,
  xlabel={配体浓度（以 $K_d$ 为单位）}, ylabel={受体占有率},
  xlabel style={font=\small, at={(axis description cs:0.5,-0.13)}, anchor=north},
  ylabel style={font=\small},
  tick label style={font=\small}, xtick={0.01,0.1,1,10,100}, ytick={0,0.5,1},
  clip=false,
]
\addplot[very thick, omDef, domain=0.01:100, samples=100] {x/(1+x)};
\draw[dashed, gray] (axis cs:1,0) -- (axis cs:1,0.5) -- (axis cs:0.01,0.5);
\node[font=\scriptsize, anchor=north west] at (axis cs:1.1,0.48) {$[L] = K_d$：半数被占据};
\node[font=\scriptsize, anchor=south east] at (axis cs:0.1,0.1) {$K_d/10$ 时 10\%};
\node[font=\scriptsize, anchor=north west] at (axis cs:10,0.92) {$10\,K_d$ 时 90\%};
\end{axis}
\end{tikzpicture}

{\small 对数坐标下的受体占据率随配体浓度的变化，
细胞的应答从十分之一升至十分之九，
在以 $K_d$ 为中心、跨度为一百倍的浓度范围内。}
\end{omfigure}

\section{表面受体}

\begin{proposition}[G 蛋白偶联受体与环 AMP]\label{prop:b2:cell-signalling:gpcr}
\index{G 蛋白}\index{G 蛋白偶联受体}\index{环 AMP}\index{第二信使}\index{蛋白激酶 A}
最大的受体家族——人体内约 \num{800} 个基因编码此类受体，
它们结合\omterm{def:b2:cell-signalling:kinds}{激素}、\omterm{def:b2:cell-signalling:kinds}{神经递质}、气味和光，是一类
跨膜七次的蛋白质。配体在膜外结合，
使其在膜内的构象发生改变，进而激活\emph{G 蛋白}：
这是一个由三个亚基组成的“开关”，它把结合的 GDP 换成 GTP，随后解离，
在结合 GTP 的形式下，数秒内调控一种效应酶，
之后水解 GTP，使自身关闭。经典的效应酶
是\emph{腺苷酸环化酶}，它以 ATP 为底物生成\emph{环 AMP}，
这是一种小分子、可扩散的\emph{\omterm{prop:b2:cell-signalling:gpcr}{第二信使}}，其浓度在数秒内
由 \qty{0.1}{\micro mol/L} 升至数微摩尔，
并激活\emph{蛋白激酶 A}，它是一种酶，
可使数十种靶蛋白磷酸化，从而将其开启或关闭，
信号由磷酸二酯酶终止，该酶破坏环 AMP，
磷酸酶则去除磷酸基团。
肾上腺素作用于肝细胞时即按此途径进行：受体 $\to$ G 蛋白
$\to$ 环化酶 $\to$ 环 AMP $\to$ 激酶 A $\to$ 磷酸化酶激酶
$\to$ 糖原磷酸化酶 $\to$ 葡萄糖。\omterm{prop:b2:cell-signalling:glucose}{胰高血糖素}沿同一途径发挥作用；
其他 G 蛋白则抑制环化酶，或激活磷脂酶 C，
后者释放另外两种\omterm{prop:b2:cell-signalling:gpcr}{第二信使}：
\emph{三磷酸肌醇}——它打开钙储存库——以及
二酰甘油。
\end{proposition}
\begin{proof}[证据]
萨瑟兰（1957 年）发现，向肝匀浆中加入肾上腺素后，
只有存在膜时才释放葡萄糖；膜能产生一种耐热因子，
该因子单独即可激活无膜组分中的磷酸化酶，
并且该因子是环 AMP——第一个\omterm{prop:b2:cell-signalling:gpcr}{第二信使}，
也是首次证明\omterm{def:b2:cell-signalling:kinds}{激素}经胞内中间体，
在细胞表面发挥作用的例证。
罗德贝尔和吉尔曼（20 世纪 70 年代）发现，环化酶需要 GTP，
并分离出在受体与酶之间传递信号的 G 蛋白；
霍乱毒素可使 G 蛋白锁定在“开”的状态，
导致肠道每天分泌数升液体。
这一疾病正是信号传导的缺陷。
\end{proof}

\begin{omfigure}
\begin{tikzpicture}[every node/.style={font=\scriptsize}, >=stealth]
  % membrane
  \fill[omExo!30] (0,2.0) rectangle (12,2.5);
  \node[anchor=east] at (-0.1,2.25) {膜};
  % receptor
  \draw[thick, fill=omDef!40] (1.0,1.7) rectangle (1.8,2.8); \node[anchor=east] at (0.95,2.65) {受体};
  \fill[omProp] (1.4,3.35) circle (0.12); \node[omProp, above] at (1.4,3.45) {肾上腺素};
  \draw[->, omProp] (1.4,3.23) -- (1.4,2.85);
  % G protein
  \draw[thick, fill=omDef!25] (2.6,1.5) ellipse (0.72 and 0.36); \node at (2.6,1.5) {G 蛋白};
  \node[below] at (2.6,1.1) {GDP $\to$ GTP};
  \draw[->, thick] (1.8,2.0) -- (2.2,1.75);
  % cyclase
  \draw[thick, fill=omDef!40] (4.2,1.7) rectangle (5.2,2.8); \node[above] at (4.7,2.85) {腺苷酸环化酶};
  \draw[->, thick] (3.15,1.5) -- (4.2,1.85);
  % cAMP
  \node[omThm, anchor=west] at (5.35,1.55) {ATP $\to$ 环 AMP};
  \draw[->, thick, omThm] (4.7,1.7) -- (5.2,1.15);
  % PKA
  \draw[thick, fill=omThm!20] (7.3,0.9) ellipse (0.7 and 0.35); \node at (7.3,0.9) {激酶 A};
  \draw[->, thick, omThm] (6.4,0.9) -- (6.6,0.9);
  % targets
  \draw[->, thick] (8.0,0.9) -- (9.0,0.9);
  \node[anchor=west, align=left] at (9.0,0.9) {磷酸化酶激酶 $\to$ 磷酸化酶 \\$\to$ 糖原 $\to$ 葡萄糖};
  % amplification annotations
  \node[gray, align=center] at (1.4,0.3) {1 个激素};
  \node[gray, align=center] at (4.7,0.3) {$\sim 10^{2}$ 个 G 蛋白};
  \node[gray, align=center] at (7.3,0.3) {$\sim 10^{4}$ 个环 AMP};
  \node[gray, align=center] at (10.2,0.3) {$\sim 10^{6}$ 个葡萄糖};
  % off switches
  \node[omExo!70!black, align=center, anchor=south] at (9.6,2.7) {关闭：GTP 水解（数秒），\\磷酸二酯酶破坏环 AMP，\\磷酸酶去除磷酸化修饰};
\end{tikzpicture}

{\small 从肾上腺素到葡萄糖的环 AMP 途径。每一级
都会激活下一级的大量分子：一个\omterm{def:b2:cell-signalling:kinds}{激素}分子
最终可产生百万个产物分子。}
\end{omfigure}

\begin{theorem}[级联放大]\label{thm:b2:cell-signalling:amplification}
\index{级联放大}
若级联的 $n$ 个阶段中，每一阶段在自身激活期间可激活下一阶段的 $a$ 个分子，
则一个信使分子将产生
$a^{n}$ 个终产物分子。设有四个放大阶段，
每级放大倍数均为一百：受体激活一百个 G 蛋白，
环化酶生成一百个环 AMP，激酶
使一百个酶磷酸化，酶释放一百个葡萄糖分子，
则 $a^{n} = 10^{8}$；实际上约为一百万个，
在数秒内由 $10^{-9}$ mol/L 的\omterm{def:b2:cell-signalling:kinds}{激素}产生。放大作用
使\omterm{def:b2:cell-signalling:kinds}{激素}能在比其所调控代谢物低百万倍的浓度下发挥效应，
也正因如此，级联的每一阶段都必须能被关闭：
无法关闭的放大器即是一种疾病
（如霍乱，
以及许多癌症中信号蛋白被锁定在“开”的状态）。
\end{theorem}
\begin{proof}
第 1 阶段产生 $a$ 个活化分子；每个分子又在
第 2 阶段产生 $a$ 个，得 $a^{2}$；由归纳法，得 $a^{n}$，即经过 $n$ 个阶段后的结果。
每一阶段的增益，等于被激活的催化剂在失活前
所能完成的底物周转次数，即其催化速率乘以
活性持续时间。
\end{proof}

\begin{proposition}[受体酪氨酸激酶与钙]\label{prop:b2:cell-signalling:rtk}
\index{受体酪氨酸激酶}\index{胰岛素受体}\index{钙信号}\index{钙调蛋白}
第二类表面受体——包括\omterm{prop:b2:cell-signalling:rtk}{胰岛素受体}、
驱动细胞分裂的生长因子受体，
以及发育过程中的诱导信号受体——
是\emph{\omterm{prop:b2:cell-signalling:rtk}{受体酪氨酸激酶}}：单次跨膜的蛋白质。
当配体使两个此类受体靠近时，它们相互磷酸化酪氨酸残基，
从而形成对接位点，供一组衔接蛋白结合。
由此同时派生出若干条级联：一条由三个激酶组成的链
（MAP 激酶级联），最终进入细胞核，开启生长相关基因；
一种脂质激酶，其产物招募介导胰岛素代谢效应的激酶，
使葡萄糖转运蛋白转移至膜上，并开启糖原合成；
此外还有其他途径。第三种普遍存在的\omterm{prop:b2:cell-signalling:gpcr}{第二信使}是\emph{钙}：
细胞质中的钙浓度维持在 \qty{0.1}{\micro mol/L}，
比外部低一万倍。它可经信号或电压开启的通道流入细胞，
或在三磷酸肌醇作用下从内质网释放，
并在毫秒内升至微摩尔浓度。
与\emph{\omterm{prop:b2:cell-signalling:rtk}{钙调蛋白}}结合后，钙激活激酶和泵；
与肌钙蛋白结合后，则触发收缩
（见\cref{ch:b2:muscle-movement}）。
泵在一秒内将钙送回储存库。
无论受体为何种，细胞内部都使用若干种“语言”，
——环 AMP、钙和磷酸化——
而特异性在于每种细胞类型拥有哪些蛋白质来响应这些信号。
\end{proposition}

\begin{omfigure}
\includegraphics[width=0.6\linewidth]{images/book4/ai/b4-19-adrenaline-runner.jpg}

{\small 级联发挥作用。受惊时肾上腺髓质释放肾上腺素；
数秒之内心跳加速，肝脏大量释放葡萄糖，
气道扩张——一种信使，
多种受体，每个细胞内实现百万倍放大。}
\end{omfigure}

\section{细胞内的受体}

\begin{proposition}[核受体]\label{prop:b2:cell-signalling:nuclear}
\index{核受体}\index{类固醇激素}
类固醇、甲状腺素、视黄酸和维生素 D 可穿过膜，
与位于细胞质或细胞核中的受体结合，这些受体
本身即是\emph{\omterm{prop:b2:cell-differentiation:transcriptional}{转录因子}}：\omterm{def:b2:cell-signalling:kinds}{激素}–受体复合物
结合特定的 DNA 序列，并使其靶基因
开启或关闭。此类反应较慢——需数分钟至数小时才能合成
mRNA 和蛋白质——但持续时间长：皮质醇在数小时内诱导
糖异生酶，雌二醇数天内促使\omterm{def:b2:reproductive-hormones:cycle}{子宫内膜}增厚，
睾酮数周内使肌肉增长，甲状腺素则决定每个细胞的
代谢率。部分类固醇也可经细胞表面受体快速起效，
部分表面通路还能抵达细胞核（如 MAP 激酶级联），
因此二者之间的界限并非绝对，
但有一条规律成立：脂溶性信使
通过改变基因来改变细胞\emph{是什么}，
而水溶性信使通过改变酶来改变细胞
\emph{做什么}。
\end{proposition}

\begin{omfigure}
\begin{tikzpicture}[every node/.style={font=\scriptsize}, >=stealth]
  \foreach \x/\t in {0/{水溶性：肾上腺素、胰岛素、胰高血糖素}, 7.4/{脂溶性：皮质醇、雌二醇、甲状腺素}} \node[anchor=south, font=\scriptsize\bfseries] at (\x+2.6,3.4) {\t};
  % left: surface receptor
  \fill[omExo!30] (0,2.4) rectangle (5.2,2.8);
  \draw[thick, fill=omDef!40] (1.6,2.1) rectangle (2.3,3.1);
  \fill[omProp] (1.95,3.35) circle (0.1);
  \draw[->, thick] (1.95,2.1) -- (1.95,1.6) node[below, align=center] {第二信使、\\激酶：在数秒内 \\切换的酶};
  \node[align=center] at (4.5,1.0) {改变细胞 \\的行为};
  % right: nuclear receptor
  \begin{scope}[xshift=7.4cm]
    \fill[omExo!30] (0,2.4) rectangle (5.2,2.8);
    \fill[omProp] (1.5,3.35) circle (0.1);
    \draw[->, omProp] (1.5,3.25) -- (1.5,1.9);
    \draw[thick, fill=omThm!20] (1.5,1.5) ellipse (0.72 and 0.3); \node at (1.5,1.5) {受体};
    \draw[thick] (0.6,0.4) -- (4.6,0.4); \draw[thick] (0.6,0.25) -- (4.6,0.25); \node[below] at (2.6,0.2) {DNA：靶基因开启或关闭，历时数小时};
    \draw[->, thick] (1.5,1.2) -- (2.4,0.5);
    \node[align=center] at (4.0,1.4) {改变细胞 \\的性质};
  \end{scope}
\end{tikzpicture}

{\small 两种设计。水溶性信使作用于细胞表面，
数秒内改变酶的活性；脂溶性信使则进入细胞，
与作为\omterm{prop:b2:cell-differentiation:transcriptional}{转录因子}的受体结合，并在数小时内改变
基因表达。}
\end{omfigure}

\section{实例：胰岛素、胰高血糖素与血糖调定点}

\begin{proposition}[两种激素与一个调定点]\label{prop:b2:cell-signalling:glucose}
\index{胰岛素}\index{胰高血糖素}\index{血糖}
血糖被维持在约 \qty{5}{mmol/L}（\qty{0.9}{g/L}）附近，
由胰岛的两种拮抗\omterm{def:b2:cell-signalling:kinds}{激素}共同完成。进食后，
血糖升高被 $\beta$ 细胞感知——它们
代谢葡萄糖，随 ATP 升高而关闭钾通道，
发生去极化，允许钙内流并释放\emph{胰岛素}——胰岛素
经其酪氨酸激酶受体，使肌肉和脂肪摄取
葡萄糖（将转运蛋白转移至膜上），肝脏
将其以糖原和脂肪的形式储存，所有组织停止分解脂肪。
两餐之间，血糖下降使 $\alpha$ 细胞释放
\emph{\omterm{prop:b2:cell-signalling:glucose}{胰高血糖素}}，它通过肝脏中的环 AMP 分解糖原，
并利用氨基酸合成新的葡萄糖；肾上腺素在紧急情况下
作用相同但更快，皮质醇和生长\omterm{def:b2:cell-signalling:kinds}{激素}则在数小时内
发挥作用。每种\omterm{def:b2:cell-signalling:kinds}{激素}的分泌量与其所纠正的偏差成正比，
即一个具有两个效应器的负反馈，两者各负责
一个方向的调节，并将调定点维持在其五分之一以内。大脑
每天消耗 \qty{120}{g} 葡萄糖且不储存，是
该调节回路所保护的器官；1 型糖尿病即
$\beta$ 细胞丧失，其未经治疗的病程——
血糖达 \qty{20}{mmol/L}，因缺乏胰岛素而分解脂肪，\omterm{def:b2:blood-circulation:blood}{血液}中出现酸——
正是去掉一条反馈支路后的调节回路表现。
\end{proposition}

\begin{omfigure}
\begin{tikzpicture}
\begin{axis}[
  width=0.82\linewidth, height=5.6cm,
  axis x line=bottom, axis y line=left,
  xmin=0, xmax=6, ymin=0, ymax=1.1,
  xlabel={进食后的小时数}, ylabel={相对于峰值},
  xlabel style={font=\small, at={(axis description cs:0.5,-0.12)}, anchor=north},
  ylabel style={font=\small},
  tick label style={font=\small}, xtick={0,1,2,3,4,5,6}, ytick={0,0.5,1},
  legend style={font=\scriptsize, at={(0.5,-0.28)}, anchor=north, draw=none, legend columns=3},
  clip=false,
]
\addplot[very thick, omDef, domain=0:6, samples=80] {0.55 + 0.45*exp(-((x-0.9)^2)/0.5) - 0.08*exp(-((x-4)^2)/2)};
\addlegendentry{血糖（0.55 处为 5 mmol/L）}
\addplot[very thick, omThm, domain=0:6, samples=80] {0.15 + 0.85*exp(-((x-1.1)^2)/0.45)};
\addlegendentry{胰岛素}
\addplot[very thick, omProp, domain=0:6, samples=80] {0.75 - 0.5*exp(-((x-1.2)^2)/0.8) + 0.25*(1-exp(-x/5))*0.6};
\addlegendentry{胰高血糖素}
\end{axis}
\end{tikzpicture}

{\small 进食后的血糖、胰岛素与\omterm{prop:b2:cell-signalling:glucose}{胰高血糖素}。胰岛素随血糖升高而升高，
并使其下降；随血糖降向调定点，
\omterm{prop:b2:cell-signalling:glucose}{胰高血糖素}升高以将其维持于此。}
\end{omfigure}

\begin{example}[两种系统的速度、作用范围与代价]\label{ex:b2:cell-signalling:compare}
\omterm{def:b2:blood-circulation:blood}{血液}中的肾上腺素约一分钟内到达每个细胞
（循环时间），作用持续数分钟；神经冲动在毫秒内到达靶点，
作用持续毫秒量级。\omterm{def:b2:cell-signalling:kinds}{激素}对所及的每个细胞几乎没有代价，
一纳摩尔分子散布于全身，但无法定向作用于单个细胞；
神经则能精确作用于一个细胞，
却必须建立并维持一条通往它的导线。机体两者兼用：
\cref{ch:b2:blood-pressure}中的交感神经在一秒之内作用于
\omterm{def:b2:heart:anatomy}{心脏}和\omterm{def:b2:blood-circulation:vessels}{小动脉}，而肾上腺髓质
本身是经过改造的交感神经节，
将同一种分子释入\omterm{def:b2:blood-circulation:blood}{血液}，为全身提供补充。
在突触处两种设计相遇：\omterm{def:b2:cell-signalling:kinds}{神经递质}是一种信使，
通过与受体结合发挥作用，机制相同——开放的通道，
切换的 G 蛋白——只是作用距离由
一米缩小到二十纳米。
\end{example}

\section{练习}

\begin{exercise}[$\star$]\label{exo:b2:cell-signalling:1}
将下列物质分类为内分泌、旁分泌、自分泌或突触信号，并判断
其为水溶性还是脂溶性：胰岛素、皮质醇、突触处的乙酰胆碱、
伤口处的组胺、雌二醇、肾上腺分泌的肾上腺素、
内皮分泌的一氧化氮。
\end{exercise}

\begin{exercise}[$\star$]\label{exo:b2:cell-signalling:2}
列出从肝细胞膜上的肾上腺素到\omterm{def:b2:blood-circulation:blood}{血液}中葡萄糖的各个步骤，
并说出每一步的关闭机制。
\end{exercise}

\begin{exercise}[$\star$]\label{exo:b2:cell-signalling:3}
比较表面受体与\omterm{prop:b2:cell-signalling:nuclear}{核受体}：它们与什么结合，
位于何处，作用速度多快，改变什么。
\end{exercise}

\begin{exercise}[$\star$]\label{exo:b2:cell-signalling:4}
本章中提到的三种通用\omterm{prop:b2:cell-signalling:gpcr}{第二信使}是什么，
各自如何升高，又如何被清除？
\end{exercise}

\begin{exercise}[$\star\star$]\label{exo:b2:cell-signalling:5}
某受体的 $K_d = \qty{2}{nmol/L}$。计算其在
0.2、2、20 和 \qty{200}{nmol/L} 时的占有率。若细胞的受体数目加倍，
则变化的是占有率，还是被占据受体的数目？
\end{exercise}

\begin{exercise}[$\star\star$]\label{exo:b2:cell-signalling:6}
某级联包含四个阶段，增益分别为 50、200、100 和 300。计算
总放大倍数。若\omterm{def:b2:cell-signalling:kinds}{激素}浓度为 \qty{1}{nmol/L}，且细胞中
\qty{1}{\%} 的受体（每个细胞 $10^{4}$ 个）被占据，
则细胞每秒产生多少个产物分子，假设每一阶段的增益
均以每秒计？
\end{exercise}

\begin{exercise}[$\star\star$]\label{exo:b2:cell-signalling:7}
霍乱毒素使肠道细胞中的刺激性 G 蛋白锁定在结合 GTP 的
状态。请逐步解释患者为何每天失去数升液体，
以及为何在细菌被清除之后，该效应仍会持续
数日。
\end{exercise}

\begin{exercise}[$\star\star$]\label{exo:b2:cell-signalling:8}
当某通道开放时，胞质钙浓度由 \qty{0.1}{} 升至 \qty{1}{\micro
mol/L}，细胞体积为 \qty{2000}{\micro m^3}。进入的
钙离子有多少个？若单个开放通道每秒通过 $10^{6}$
个离子，则一个通道需要多长时间，并说明细胞为何
会同时开放多个通道？
\end{exercise}

\begin{exercise}[$\star\star$]\label{exo:b2:cell-signalling:9}
胰岛素和\omterm{prop:b2:cell-signalling:glucose}{胰高血糖素}都作用于肝脏，一种经酪氨酸
激酶，另一种经 cAMP。请解释肝细胞如何整合这两种信号，
两者同时升高时糖原会发生什么变化，
以及为何二者的比值比任一浓度水平都更重要。
\end{exercise}

\begin{exercise}[$\star\star\star$]\label{exo:b2:cell-signalling:10}
肾上腺素在 \qty{5}{s} 内使肝细胞中的 cAMP 升高，并在
\qty{30}{s} 内使葡萄糖输出增加；皮质醇则使肝脏糖异生酶
在 \qty{4}{h} 内升高。利用扩散、转录（速率为每秒 \qty{50}{}
个核苷酸，基因长 \qty{3000}{} 个核苷酸）、翻译（速率为每秒
\qty{5}{} 个氨基酸，蛋白质长 \num{500} 个残基）以及
积累数千个酶分子的需要，说明
这两个时间尺度。
\end{exercise}

\begin{exercise}[$\star\star\star$]\label{exo:b2:cell-signalling:11}
一个\omterm{def:b2:genome-diversification:mutation}{突变}使某生长因子\omterm{prop:b2:cell-signalling:rtk}{受体酪氨酸激酶}在无配体时发生二聚化，
从而发出信号。请预测该细胞的行为，
解释为何这是癌症中常见的一步，并说明
一种阻断该激酶 ATP 结合位点的药物会产生什么作用。
\end{exercise}

\begin{exercise}[$\star\star\star$]\label{exo:b2:cell-signalling:12}
“内分泌系统是广播式的；神经系统是布线式的。”
请从速度、作用范围、特异性和代价等方面讨论两种设计，
并说明两者各自在何处不可或缺，以及机体在何处对同一信息
同时使用两种方式。
\end{exercise}

\section{问题：从血液到基因的一种激素}

\begin{problem}[{周末问题——追踪肾上腺素从肾上腺到肝细胞葡萄糖，
并计算其浓度、
受体占有率、级联增益与时间，再追踪皮质醇到基因，
并分析血糖回路，
最后落在占有率、放大倍数以及每种激素所需的时间上}]
\label{pb:b2:cell-signalling:1}
已知数据：遇到惊吓时，肾上腺髓质释放 \qty{1}{nmol} 肾上腺素，
进入 \qty{5}{L} \omterm{def:b2:blood-circulation:blood}{血液}；肾上腺素在肝脏受体上的 $K_d$ 为
\qty{1}{nmol/L}；一个肝细胞（体积 \qty{5000}{\micro m^3}）
携带 $2\times 10^{4}$ 个受体。每秒激活的级联增益为：受体 $\to$ G 蛋白 20，
G 蛋白 $\to$ cAMP 100
（环化酶在 G 蛋白 3 秒寿命内的周转），
cAMP $\to$ 激酶 A 1（化学计量比），激酶 A $\to$ 磷酸化酶
激酶 50，$\to$ 磷酸化酶 50，$\to$ 葡萄糖 1000。肝脏储存
\qty{100}{g} 糖原。皮质醇：转录速率为
每秒 \qty{50}{} 个核苷酸，基因长 \qty{4000}{} 个核苷酸，
翻译速率为每秒 \qty{5}{} 个残基，酶由 \num{600} 个残基组成，
每个 mRNA 对应 \num{20} 个核糖体，每小时制成 \num{100}
个 mRNA，效应需要 $10^{5}$ 个酶分子。

\medskip
\noindent\textbf{第一部分——\omterm{def:b2:cell-signalling:kinds}{激素}抵达。}
\begin{enumerate}
  \item 忽略清除作用，计算释放后
        肾上腺素的\omterm{def:b2:blood-circulation:blood}{血浆}浓度。
  \item 计算\omterm{thm:b2:cell-signalling:occupancy}{受体占有率}，以及
        每个肝细胞中被占据的受体数。
  \item 肾上腺素的清除半衰期为 \qty{2}{min}，
        经过 \qty{6}{min} 后，其浓度与占有率各为多少？
  \item 肾上腺素从肾上腺到达肝脏
        需要多长时间？设循环时间约为一分钟。为何
        通向肝脏的交感神经作用更快？
  \item 若一个细胞有 $2\times 10^{5}$ 个受体，且 $K_d$ 相同，
        会发生什么变化？
  \item 为何 $K_d$ 须与循环中的浓度相匹配 ---
        若 $K_d = \qty{1}{\micro mol/L}$ 的受体
        遇到纳摩尔级的\omterm{def:b2:cell-signalling:kinds}{激素}，会怎样？
\end{enumerate}

\medskip
\noindent\textbf{第二部分——级联反应。}
\begin{enumerate}[resume]
  \item 计算从一个被占据的受体到每秒葡萄糖分子
        数的总增益。
  \item 计算细胞在第一
        秒内制成的 cAMP 分子数，以及它们在细胞体积中达到的浓度。
  \item 计算一个肝细胞在
        第 2 题的占有率下，每秒及每分钟的葡萄糖输出量。
  \item 肝脏有 $2\times 10^{11}$ 个细胞，计算各级增益所预测的
        输出量（以克每分钟计），并与
        \qty{5}{g/min} 的实际观测值相比较。是什么限制了实际
        输出？
  \item 若以此速率计算，\qty{100}{g} 糖原可维持
        多久？为何受惊时它不会耗尽？
  \item 指出每个放大环节上的关闭开关，并说明
        哪一个作用最快。
  \item 若存在磷酸二酯酶抑制剂（咖啡因），预测
        其对级联反应及葡萄糖输出的影响。
\end{enumerate}

\medskip
\noindent\textbf{第三部分——从皮质醇到基因。}
\begin{enumerate}[resume]
  \item 计算转录一个 mRNA 所需的时间。
  \item 计算翻译一个酶分子所需的时间，以及
        每小时每个 mRNA 制成的酶分子数（设
        每个 mRNA 存活一小时，核糖体依次起始）。
  \item 计算每小时由
        \num{100} 个 mRNA 制成的酶分子数，以及达到 $10^{5}$ 所需的时间。
  \item 将皮质醇进入、结合并到达 DNA 所需的时间
        （以分钟计）相加，并解释观察到的 \qty{4}{h}。
  \item 为何机体不对皮质醇的任务使用
        类似肾上腺素的级联，而对受惊反应
        又为何不采用皮质醇的机制？
  \item 皮质醇还会增加肝细胞中的
        肾上腺素受体数量。试解释慢速\omterm{def:b2:cell-signalling:kinds}{激素}如何放大
        快速\omterm{def:b2:cell-signalling:kinds}{激素}的作用，以及这对快速\omterm{def:b2:cell-signalling:kinds}{激素}的占有率曲线有何影响。
\end{enumerate}

\medskip
\noindent\textbf{第四部分——反馈环路。}
\begin{enumerate}[resume]
  \item 进食后葡萄糖由 5 升至 \qty{8}{mmol/L}。
        描述 $\beta$ 细胞由葡萄糖到释放胰岛素的反应，
        并指出其中的偶联。
  \item 胰岛素通过酪氨酸激酶起作用，\omterm{prop:b2:cell-signalling:glucose}{胰高血糖素}则通过
        cAMP 起作用。解释单个细胞为何能同时响应两者，以及
        肝脏中的糖原酶如何由二者的比例所决定。
  \item 某患者没有 $\beta$ 细胞。预测其空腹及
        餐后血糖，并解释为何脂肪被分解且
        并产生酸。
  \item 注射的胰岛素没有反馈作用。解释
        通常剂量下作用异常持续所带来的危险。
  \item 比较葡萄糖环路与
        \cref{ch:b2:blood-pressure}中的\omterm{def:b2:blood-pressure:baroreflex}{压力反射}：感受器、效应器、增益与速度。
  \item 陈述结果：释放后肾上腺素的占有率、
        级联增益、肝脏的葡萄糖输出量，以及
        皮质醇所需的时间。
\end{enumerate}
\end{problem}
